% !TEX root = ../metatime_monograph.tex
\chapter{Separable Universal Candidate Family}
\label{app:separable_candidate_family}


% CITATION_CONTEXT_V10_9_BEGIN
\paragraph{Established context.}
The candidate-family freeze follows falsifiability, preregistration, and
multiplicity-aware interpretation \cite{Popper1959,Nosek2018,Munafo2017}.  Its
geometric and arithmetic ingredients are externally contextualized by Berry,
Collatz, and Ramanujan literature
\cite{berry1984,Lagarias1985,Tao2022Collatz,Ramanujan1918}, while the three
specific maps and their Yukawa-ratio predictions are original TIR candidates.
% CITATION_CONTEXT_V10_9_END

\section{Scope and status}

Appendix~\ref{app:observable_identifiability} separated a
sector-invariant baseline from a generation-varying release.  The present
appendix freezes exactly three algebraic candidates before any qualifying
future Higgs-coupling likelihood is inspected.

The status is

\[
\boxed{
\begin{gathered}
\text{candidate family frozen}
\;\land\;
\text{no mass benchmark}\\
\text{no candidate selected}
\;\land\;
\text{canonical promotion denied}
\end{gathered}
}
\]

\section{Structural coordinates}

Let

\[
S_f=A_f^{\rm rep}+C_{\rm chir}
\]

be the generation-invariant charged-sector action.  The corrected v1.0
representation table and the universal charged-Dirac chiral cost give

\[
S_\ell=4.5,
\qquad
S_d=3.0876164691516155,
\qquad
S_u=2.920949802484949.
\]

The v1.0 generation damping coordinates are

\[
G_1=0.8471633855025916,
\qquad
G_2=0.6088811663290957,
\qquad
G_3=0.23926393244694075.
\]

The mandatory v2.1 Ramanujan release units are

\[
R_1=0.4754443238393299,
\qquad
R_2=1.0262897039931138,
\qquad
R_3=2.442442115365772.
\]

No observed fermion mass enters these coordinates.

\section{Separable architecture}

Every candidate has the form

\begin{equation}
\ln y_{f,g}=F(S_f)+D(G_g,R_g).
\label{eq:v107_separable_architecture}
\end{equation}

The same functions $F$ and $D$ act on charged leptons, down-type quarks, and
up-type quarks.  This makes the sector and generation coordinates
experimentally separable.

\section{Candidate C1: linear action}

\begin{equation}
F_1(S)=-S,
\qquad
D_1(G,R)=-G+R.
\label{eq:v107_candidate_c1}
\end{equation}

This is the direct exponential-damping interpretation of the v1.0 structural
action.

\section{Candidate C2: Collatz quarter-power action}

\begin{equation}
F_2(S)=-S^{3/4},
\qquad
D_2(G,R)=-G^{3/4}+R.
\label{eq:v107_candidate_c2}
\end{equation}

The exponent is not refitted.  It is the accelerated-Collatz quarter-power
already frozen in Appendix~\ref{app:collatz_quarter_power}.

\section{Candidate C3: inverse quarter-power information potential}

\begin{equation}
F_3(S)=\frac{S^{-3/4}}{L_3\kappa},
\qquad
D_3(G,R)=\frac{G^{-3/4}}{L_3\kappa}+R.
\label{eq:v107_candidate_c3}
\end{equation}

This is the separable universal-potential form corresponding to the inverse
action coordinate used in the original quarter-power bridge.  Its large
predicted generation range is retained rather than removed after inspection.

\section{Orthogonal prospective observables}

Charm and the muon are both second-generation states.  Therefore

\begin{equation}
\ln\frac{y_c}{y_\mu}=F(S_u)-F(S_\ell),
\label{eq:v107_sector_observable}
\end{equation}

because $D(G_2,R_2)$ cancels exactly.  The primary class-A observable is thus
the first qualifying post-29 July 2026 joint ATLAS/CMS direct charm-to-muon
Higgs-coupling likelihood.

Charm and top share the up-type/heavy sector in the frozen architecture, so

\begin{equation}
\ln\frac{y_c}{y_t}=D(G_2,R_2)-D(G_3,R_3),
\label{eq:v107_generation_observable}
\end{equation}

and the sector baseline cancels.  The primary class-B observable remains the
first qualifying post-29 July 2026 joint ATLAS/CMS direct charm-to-top
Higgs-coupling likelihood.

The previous use of $y_c/y_\tau$ as the primary class-A test is superseded: it
crosses both sector and generation and therefore does not isolate the sector
baseline.

\section{Frozen predictions}

\begin{center}
\begin{tabular}{lrr}
\toprule
Candidate & $y_c/y_\mu$ & $y_c/y_t$\\
\midrule
C1 linear action & 4.850346751338371 & 0.16766796647328305\\
C2 quarter-power action & 2.3521800134268784 & 0.17147213462587316\\
C3 inverse quarter potential & 6.858021228826222 & $2.8101955040512466\times10^{-11}$\\
\bottomrule
\end{tabular}
\end{center}

All three candidates remain active.  No current mass table or already visible
experimental likelihood is used to select among them.

\section{Cross-transfer and decision rule}

The same maps emit predictions for $y_s/y_\mu$, $y_u/y_e$, $y_\mu/y_e$,
$y_\tau/y_\mu$, $y_c/y_u$, and $y_t/y_c$.  These are frozen diagnostics, not
selection variables.

For each candidate, PASS requires the frozen ratio to lie inside the published
95\% confidence region of the assigned first qualifying likelihood.  Failure is
retained without changing the formula or substituting another observable.  The
three-member family requires a multiplicity-aware interpretation.

The executable ledger is generated by

\path{TIR/validation/separable_universal_candidate_family_v10_7.py}

and carries fingerprint

\path{2bbf9985b6a4fc9d19d039e117b15eb446a51a3e8828a2894c8d797bb35f23f4}.
